\scanpage{160}
\mspage{105}
pour $X$ dans $\deleted{\unsure{C}}\mathcal M_{/S}$, $C(X)$ représente le \inserted{\unsure{contra.}}foncteur suivant dans $\mathcal M\deleted{_{/S}}$ : pour tout $T$ dans $\mathcal M$,
\[
\begin{aligned}
C(X)(T)=\{(x,c)\mid x\in X(T),\quad &\unsure{c}:f_X(x)\xrightarrow{c}\unsure{s}\ \text{un}\\
&\text{chemin dans }S(T),\\
&\text{de source } f_X(x)\}
\end{aligned}
\]
où $f_X:X\longrightarrow S$ est le morphisme structural. Les ``calculs'' \unsure{faits} il y a deux (?) jours montrent qu'on a de $\eqsim$ une interprétation
\[
\begin{aligned}
C(C(X))(T)=\{(x,c,c')\mid x\in X(T),\quad &f_X(x)\mathrel{\overset{c}{\rule{1.6em}{.4pt}}}s\xrightarrow{c'}\unsure{t}\deleted{\unsure{\}}}\\
&c,\text{ un chemin de source }f_X(x),\ c'\\
&\text{un chemin de source égale}\\
&\text{à }\beta(c)\ \}
\end{aligned}
\]
et les applications $i_{C(X)}$ et $\deleted{\unsure{e}}C(i_X)$ s'identifient respectivement à
\[
\begin{aligned}
(x,c)&\longmapsto(x,1_{f(x)},c)\deleted{\unsure{f(x)}}\inserted{\unsure{\text{i.e.}|\mathrm{id}}} && f(x)\xrightarrow{1_x}f(x)\xrightarrow{c}s\\
\deleted{\unsure{\&}}\ (x,c)&\longmapsto(x,c,1_{\beta(c)}) && f(x)\longrightarrow s\xrightarrow{1_s}s\text{.}
\end{aligned}
\]
Pour les comparer, on a une envie irrésistible d'écrire
\[
c\circ1_x=1_s\circ c=c\text{,}
\]
mais en a-t-on le droit ? Il faudrait pouvoir composer les chemins, et c'est très, très difficile !

Mais on peut toujours axiomatiser. On suppose donc que pour $\deleted{\unsure{C}}X$ dans $\mathcal M$, $\mathrm{Ch}(X)$
% archive sheet number: 159
